\documentclass[tikz,border=8pt]{standalone}
\usepackage{ctex}
\usepackage{amsmath}
\usetikzlibrary{arrows.meta}
\begin{document}
\begin{tikzpicture}[>=Stealth]
  % ================= 左panel：几何截面 =================
  % 内球（实心导体）
  \filldraw[fill=gray!35, draw=black, thick] (0,0) circle (1.1);
  \node at (0.35,0.25) {$\oplus$};
  \node at (-0.4,0.35) {$\oplus$};
  \node at (0.05,-0.45) {$\oplus$};
  \node at (-0.35,-0.2) {$\oplus$};
  \node[below=2pt] at (0,-1.12) {内球 $+q$};
  % 外壳（薄导体壳）
  \draw[very thick] (0,0) circle (2.2);
  \node at (90:2.2) {$\ominus$};
  \node at (150:2.2) {$\ominus$};
  \node at (210:2.2) {$\ominus$};
  \node at (270:2.2) {$\ominus$};
  \node at (330:2.2) {$\ominus$};
  \node[above=2pt] at (0,2.22) {外壳 $-q$};
  % 间隙电场箭头（沿径向向外）
  \foreach \ang in {60,120,185,235,300}
    \draw[->,teal,thick] (\ang:1.25) -- (\ang:2.02);
  \node[teal] at (30:1.62) {$\vec{E}$};
  % 水平虚线半径 + 三个区域标记
  \draw[dashed,gray] (0,0) -- (2.75,0);
  \fill (0.55,0) circle (1.6pt);
  \node[below] at (0.55,-0.06) {\footnotesize\textcircled{\scriptsize 1}};
  \fill (1.65,0) circle (1.6pt);
  \node[below] at (1.65,-0.06) {\footnotesize\textcircled{\scriptsize 2}};
  \fill (2.75,0) circle (1.6pt);
  \node[below] at (2.75,-0.06) {\footnotesize\textcircled{\scriptsize 3}};
  \node[below=8pt] at (1.1,0) {$r_a$};
  \node[above right=2pt] at (2.2,0.02) {$r_b$};
  % 中心点
  \fill (0,0) circle (1.6pt);
  % 外部 V=0 注记（置于底部，避免与外壳圆圈相碰）
  \node[font=\footnotesize,align=left,anchor=west] at (-3.4,-3.05)
    {参考点：$V=0$ 取在无穷远；\textcircled{\scriptsize 3} 区（$r>r_b$）：$V\equiv 0$};

  \begin{scope}[xshift=8.2cm]
    % ================= 右panel：V(r) 剖面 =================
    \draw[->] (0,0) -- (3.45,0) node[below] {$r$};
    \draw[->] (0,0) -- (0,1.45) node[left] {$V(r)$};
    % 平台段 ①：r<r_a
    \draw[very thick,blue] (0.12,1) -- (1,1);
    % 下降段 ②：r_a<r<r_b, V/(q/4pi eps0 rb 比例) = 2/r - 1
    \draw[very thick,blue,domain=1:2,samples=80,smooth] plot (\x,{2/\x-1});
    % 零段 ③：r>r_b
    \draw[very thick,blue] (2,0) -- (3.3,0);
    % 参考虚线
    \draw[dashed,gray] (0,1) -- (1,1);
    \draw[dashed,gray] (1,0) -- (1,1);
    \node[left] at (0,1) {$V_{ab}$};
    \node[below] at (1,-0.02) {$r_a$};
    \node[below] at (2,-0.02) {$r_b$};
    % 区域标记
    \node[font=\footnotesize] at (0.55,1.14) {\textcircled{\scriptsize 1}};
    \node[font=\footnotesize] at (1.18,0.42) {\textcircled{\scriptsize 2}};
    \node[font=\footnotesize,above] at (2.65,0.04) {\textcircled{\scriptsize 3} $V=0$};
    % 公式注记（置于曲线上方空白处）
    \node[font=\footnotesize,anchor=west] at (1.5,0.62)
      {$\frac{q}{4\pi\varepsilon_0}\left(\frac{1}{r}-\frac{1}{r_b}\right)$};
  \end{scope}
\end{tikzpicture}
\end{document}
